\chapter*{Declaration}
\addcontentsline{toc}{chapter}{Declaration}
I hereby declare that the project report entitled \textbf{"HUMAN ACTIVITY Recognition"} 
is a record of the work carried out by me 
as part of the DSF project, 
under the guidance of Prof. Pisipati Radha Krishna. The implementation details and experimental results presented in this report are based on the submitted project notebook and associated dataset files.

\vspace{1.5cm}
\textbf{Place:} NIT Campus\hfill \textbf{Signature:} Hrishikesh Chandra\\
\textbf{Date:} 07/Nov/2026

\chapter*{Acknowledgement}
\addcontentsline{toc}{chapter}{Acknowledgement}
I express my sincere 
gratitude to Prof. Pisipati Radha Krishna for the 
guidance and support 
provided during this project. 
I also thank the Department CSE, NIT Warangal 
for providing the academic environment 
and resources required to complete the work. 
Finally, I thank my peers and everyone who 
contributed directly or indirectly to the 
successful completion of this project.

\chapter*{Abstract}
\addcontentsline{toc}{chapter}{Abstract}
Human Activity Recognition (HAR) is the task of identifying a person's physical activity from sensor measurements. This project develops a machine-learning workflow to classify six activities: walking, walking upstairs, walking downstairs, sitting, standing and laying. The supplied dataset contains 561 engineered features for each observation, with 7,352 training samples and 2,947 test samples. Three classifiers are trained and evaluated: K-Nearest Neighbours (KNN), Decision Tree and Random Forest. KNN uses a standardisation step within a machine-learning pipeline, while the Decision Tree and Random Forest are trained directly on the supplied feature vectors. Performance is evaluated using accuracy, macro-averaged precision, recall and F1-score, along with class-wise classification reports and confusion matrices. \emph{Random Forest obtains the strongest test performance, with 92.87\% accuracy and a macro F1-score of 92.70\%.} The trained models and supporting metadata are also exported for use by a Streamlit web application. The report documents the dataset, workflow, implementation, evaluation and application deployment evidence available from the project files.

\textbf{Keywords:} Human Activity Recognition, Machine Learning, KNN, Decision Tree, Random Forest, Sensor Features, Classification.
